% Worked-derivation fragment of Wansing's DK display calculus. Structural
% connectives expose formula positions; this is not presented as the complete
% rule surface to which Wansing's adequacy and cut theorems apply.
\calculussheet{sheet:display-k-wansing}{Worked-Derivation Fragment of Wansing's Display Calculus $\mathsf{DK}$}{Language: classical modal formulas with box primitive; structures $X,Y$ are generated from formulas, the unit $I$, structural negation $*X$, product $X\circ Y$, and modal structural operator $\bullet X$; display sequents are $X\Rightarrow Y$; the structural equivalence below is reversible; the listed formula rules are applied only when the principal formula is displayed; equality: none; scope: exactly the fragment used by Definition~\ref{def:p07-display-k-derivation}, not a complete transcription of Wansing's calculus.}
\noindent\emph{Controlling source:} \cite{Wansing1994}.
\hbrule{rule:display-k-bullet}{Modal Display}{\bullet X\Rightarrow Y}{X\Rightarrow\bullet Y}{reversible}
\hbrule{rule:display-k-product-left}{Product Display-L}{X\circ Y\Rightarrow Z}{X\Rightarrow *Y\circ Z}{reversible}
\hbrule{rule:display-k-product-right}{Product Display-R}{X\Rightarrow Y\circ Z}{X\circ *Z\Rightarrow Y}{reversible}
\hbrule{rule:display-k-box-right}{Box-R}{\bullet X\Rightarrow A}{X\Rightarrow\Box A}{\none}
\hbrule{rule:display-k-box-left}{Box-L}{A\Rightarrow X}{\Box A\Rightarrow\bullet X}{\none}
\hbrule{rule:display-k-imp-right}{Implication-R}{X\circ A\Rightarrow B}{X\Rightarrow A\to B}{\none}
\hbrule{rule:display-k-imp-left}{Implication-L}{X\Rightarrow A\qquad B\Rightarrow Y}{A\to B\Rightarrow *X\circ Y}{\none}
\hbrule{rule:display-k-id}{Identity}{}{A\Rightarrow A}{\none}
\hbrule{rule:display-k-monotonicity}{Structural Monotonicity}{X\Rightarrow Z}{X\circ Y\Rightarrow Z}{the symmetric antecedent and dual succedent forms are also available}
\hbrule{rule:display-k-contraction}{Structural Contraction}{X\circ X\Rightarrow Z}{X\Rightarrow Z}{with the dual succedent form}
\hbrule{rule:display-k-unit}{Structural Unit}{X\Rightarrow Z}{I\circ X\Rightarrow Z}{reversible deletion of the displayed unit}
